\part{実数の連続性}

\begin{abstract}
前章では，$\varepsilon$--$N$論法を用いることで，（実）数列の収束を厳密に定義した．
しかし，有理数列 $a_n = (1+1/n)^n$ が有理数の範囲に極限を持たないように，
「数列の行き先」が常にその集合の中に存在するかという問題は，極限の定義とは別の議論を必要とする．
そのときに取り上げられるものが，実数体を特徴づける性質である「実数の連続性」である．
実数体 $\mathbb{R}$ において，数列が「隙間」に落ち込むことなく収束先を見出せるのは，
我々が「実数の連続性」という強力な事実を「公理」として認めているからに他ならない．

本章ではまず，「有理数の集合には上限が存在しない場合がある」という具体例を確認し，
その上で，「実数の連続性」の公理を導入し，これを用いることで「単調収束定理」「アルキメデスの原理」などの命題を証明していく．
また，高校の数学でも扱った$\lim_{n \to \infty} 1/n = 0$ という直感的な事実を厳密に支える「アルキメデスの原理」が，
この連続性の公理からどのように導かれるのか，その論理的な構造を明らかにしよう．
\end{abstract}

\section{実数の連続性}

\subsection{「実数の連続性」の公理}

ここで，一般にある集合$A$が上に有界であっても，その上限が$A$の中に存在しないことの具体例を挙げよう．

\begin{example}{}{}

集合$A \coloneqq \{x \mid x \in \mathbb{Q},\ x>0,\ x^2<3\}$は上に有界であり，$3$は$A$の上界であるが，$\sup A$は$\mathbb{Q}$の中には存在せず，当然$A$の中にも存在しない．

ここで，$A$に上限が存在すると仮定して，それを$s$とおこう．$1 \in A$であるから$s \geqq 1 > 0$である．
また$s \in \mathbb{Q}$であり$\sqrt{3} \notin \mathbb{Q}$であるから，$s^2 \ne 3$である．
以下，$s^2 < 3$と$s^2 > 3$のいずれの場合にも矛盾が生じることを示す．

\medskip
\noindent\textbf{(i) $s^2 < 3$ のとき．}
$h \coloneqq \dfrac{3 - s^2}{2(2s+1)}$とおくと，$s \geqq 1$より$0 < h < 1$である．このとき$s_1 \coloneqq s + h$は$s_1 \in \mathbb{Q}$，$s_1 > 0$であり，$h^2 < h$より
\[
  s_1^2 = s^2 + 2sh + h^2 < s^2 + (2s+1)h = \frac{s^2+3}{2} < 3
\]
が成り立つ．よって$s_1 \in A$かつ$s_1 > s$となり，これは$s$が$A$の上界であることに矛盾する．

\medskip
\noindent\textbf{(ii) $s^2 > 3$ のとき．}
$ s_0 \coloneqq s - \dfrac{s^2-3}{2s} = \dfrac{s^2+3}{2s} $とおくと\footnote{この$s_0$は，$y= x^2 -3$の$x=s$における接線と$x$軸との交点の$x$座標である．このように$s_0$を定めるとうまくいく理由は，曲線が下に凸である性質を利用しているためである．接線が$x$軸と交わる点（$y=0$）において，接線より上にある元の曲線は必ず正の値をとるため，より小さな上界$s_0$を幾何学的に確実に見つけることができる．}，$s_0 > 0$であり，
\[
s_0^2 = s^2 - (s^2-3) + \frac{(s^2-3)^2}{4s^2} = 3 + \frac{(s^2-3)^2}{4s^2} > 3
\]
が成り立つ．したがって，任意の$x \in A$に対して$0 < x$かつ$x^2 < 3 < s_0^2$より$x < s_0$となるから，$s_0$は$A$の上界である．しかし，
\[
s_0 - s = - \frac{s^2-3}{2s} < 0
\]
であるから$s_0 < s$であり，これは$s$が$A$の上限（最小の上界）であることに矛盾する．

\medskip
以上より，いずれの場合も矛盾が生じるから，$A$の上限は$\mathbb{Q}$の中には存在せず，したがって$A$の中にも存在しない．
\end{example}



このような状況を踏まえ，われわれは集合$\mathbb{R}$において以下の事実を公理\footnote{数学において，最初の前提になる事実は公理と呼ばれる．公理は推論のための前提にすぎす，誰もが認める自明の事柄などではない．}として認めることにする．
\begin{axiom}{実数の連続性}{実数の連続性} \index{じっすうのれんぞくせい@実数の連続性}
	$A \subset \mathbb{R}$とする．$A \ne \varnothing$が上に有界であるとき，$\sup A$が$\mathbb{R}$の中に存在する．
\end{axiom}


ここまでで実数の連続性について確認した．以降でそれを用いて様々な命題を証明していく．


\phantomsection
\subsection{単調収束定理}
\index{たんちょうしゅうそくていり@単調収束定理}

\begin{definition}{単調増加・単調減少}{}
  実数列 $(a_n)_{n \in \mathbb{N}}$ において，
  \begin{itemize}
    \item 任意の $n \in \mathbb{N}$ に対して $a_n \leqq a_{n+1}$ が成り立つとき，$(a_n)$ は\textbf{単調増加}であるという．特に $a_n < a_{n+1}$ が成り立つときは\textbf{狭義単調増加}であるという．
    \index{たんちょうぞうかすうれつ@単調増加数列}
    \item 任意の $n \in \mathbb{N}$ に対して $a_n \geqq a_{n+1}$ が成り立つとき，$(a_n)$ は\textbf{単調減少}であるという．特に $a_n > a_{n+1}$ が成り立つときは\textbf{狭義単調減少}であるという．
    \index{たんちょうげんしょうすうれつ@単調減少数列}
  \end{itemize}
\end{definition}

\begin{theorem}{単調収束定理}{}
  \begin{enumerate}
    \item 上に有界な単調増加数列 $(a_n)_{n \in \mathbb{N}}$ は，上限 $s = \sup \{ a_n \mid n \in \mathbb{N} \}$ に収束する．
    \item 下に有界な単調減少数列 $(a_n)_{n \in \mathbb{N}}$ は，下限 $i = \inf \{ a_n \mid n \in \mathbb{N} \}$ に収束する．
  \end{enumerate}
\end{theorem}



\begin{leftbar}
  \begin{proof}
    前半の主張を示す（後半も同様に示せる）．
    数列の項からなる集合を $A \coloneqq \{a_n \mid n \in \mathbb{N}\}$ とおく．
    $A$ は空でなく，かつ仮定により上に有界であるから，実数の連続性（上限性質）により上限 $s \coloneqq \sup A$ が実数として存在する．

    上限の定義より，任意の $n \in \mathbb{N}$ に対して $a_n \leqq s$ が成り立つ．
    次に，任意の $\varepsilon > 0$ をとる．$s - \varepsilon < s$ であるから，$s - \varepsilon$ は $A$ の上界ではない．
    したがって，ある $n_0 \in \mathbb{N}$ が存在して $s - \varepsilon < a_{n_0}$ が成り立つ．

    ここで，$(a_n)$ は単調増加であるから，すべての $n \geqq n_0$ に対して $a_{n_0} \leqq a_n$ である．
    ゆえに，$n \geqq n_0$ ならば
    \[
      s - \varepsilon < a_{n_0} \leqq a_n \leqq s < s + \varepsilon
    \]
    が成り立ち，これは $\abs{a_n - s} < \varepsilon$ を意味する．
    したがって $\lim_{n \to \infty} a_n = s$ である．これが示すべきことであった．
  \end{proof}
\end{leftbar}

この事実は，$\lim_{n \to \infty} (1+1/n)^n$ のような数列の極限を求める際に非常に有用である．
実際，数列 $a_n = (1+1/n)^n$ は単調増加であることが容易に確認できる．さらに，$a_n < 3$ であることも確認できるから，単調収束定理により，$a_n$ はある実数に収束する．
この数列の極限は「ネイピア数」と呼ばれ，$e$と表される．さらに，$e$は無理数であることも知られている．

\phantomsection
\subsection{アルキメデスの原理} \index{アルキメデスのげんり@アルキメデスの原理}

\begin{theorem}{アルキメデスの原理}{}
  任意の実数 $a > 0$ および $b > 0$ に対して，
  \[
    na > b
  \]
  をみたす自然数 $n \in \mathbb{N}$ が存在する．これは，自然数の集合 $\mathbb{N}$ が $\mathbb{R}$ において上に有界ではないことと同値である．
\end{theorem}

\begin{leftbar}
  \begin{proof}
    主張の後半（$\mathbb{N}$ が上に有界でないこと）を示せば，任意の $b/a$ に対してそれより大きい $n$ を取れるため，前半も示される．
    いま，$\mathbb{N}$ が上に有界であると仮定し，矛盾を導く．
    \axref{実数の連続性} により，上限 $s \coloneqq \sup \mathbb{N}$ が実数として存在する．
    上限の定義から，$s-1$ は $\mathbb{N}$ の上界ではない．
    したがって，ある $n \in \mathbb{N}$ が存在して，$s-1 < n$ が成り立つ．
    このとき $s < n+1$ となるが，$n+1$ もまた自然数であるため，これは $s$ が $\mathbb{N}$ の上限（上界）であることに矛盾する．
    ゆえに，$\mathbb{N}$ は上に有界ではない．
  \end{proof}
\end{leftbar}

アルキメデスの原理は，実数の連続性と密接に関わっており，以下の「有理数の稠密性」と実質的に同値な概念である．

\begin{theorem}{有理数の稠密性}{} \index{QのRにおけるちゅうみつせい@$\mathbb{Q}$の$\mathbb{R}$における稠密性}
  任意の $a, b \in \mathbb{R}~(a < b)$ に対して，
  \[
    a < r < b
  \]
  をみたす有理数 $r \in \mathbb{Q}$ が存在する．
\end{theorem}

「有理数の稠密性」を仮定して「アルキメデスの原理（$\mathbb{N}$ が上に有界でないこと）」を導く過程は以下の通りである．

\begin{leftbar}
  \begin{proof}
    任意の $M > 0$ に対して，それよりも大きい自然数 $n$ が存在することを示す．
    有理数の稠密性の仮定（任意の開区間に有理数が存在すること）に基づき，開区間 $( 0, 1/M )$ を考える．
    このとき，
    \[
      0 < r < \frac{1}{M}
    \]
    を満たす有理数 $r \in \mathbb{Q}$ が存在する．
    $r$ は正の有理数であるから，自然数 $p, q$ を用いて $r = q/p$ と表せる．
    ここで，$q \geqq 1$ であるから，
    \[
      \frac{1}{p} \leqq \frac{q}{p} < \frac{1}{M}
    \]
    が成り立つ．この不等式の両辺（正数）の逆数をとることで，
    \[
      M < p
    \]
    を得る．$p$ は自然数であるから，与えられた $M$ に対してそれより大きい自然数の存在が示された．
  \end{proof}
\end{leftbar}

アルキメデスの原理と同値な命題はいくつもある．
たとえば
\[
\lim_{n \to \infty} \frac{1}{n}=0
\]
や
\[
\lim_{n \to \infty}  n = \infty
\]
などもアルキメデスの原理と同値である．

